\documentclass[11pt]{article}
\usepackage[T1]{fontenc}
\usepackage{lmodern,amsmath,amssymb,amsthm,booktabs,microtype,listings,xcolor}
\usepackage[a4paper,margin=25.5mm,top=24mm,bottom=25mm]{geometry}
\usepackage[hidelinks]{hyperref}
\usepackage{fancyhdr}
\setlength{\headheight}{14pt}
\pagestyle{fancy}\fancyhf{}\fancyhead[L]{\small Transcendence of shifted Mills minima}
\fancyhead[R]{\small Research manuscript}\fancyfoot[C]{\thepage}
\setlength{\emergencystretch}{2em}\raggedbottom
\newtheorem{theorem}{Theorem}
\newtheorem{lemma}[theorem]{Lemma}
\newtheorem{proposition}[theorem]{Proposition}
\DeclareMathOperator{\Tr}{Tr}
\DeclareMathOperator{\Nm}{N}
\definecolor{ink}{RGB}{25,49,76}
\lstset{basicstyle=\ttfamily\footnotesize,keywordstyle=\color{ink}\bfseries,
 commentstyle=\itshape\color{black!70},keepspaces=true,columns=fullflexible,
 showstringspaces=false,breaklines=true,aboveskip=8pt,belowskip=8pt}
\hypersetup{pdftitle={Transcendence of shifted Mills minima for odd multipliers up to one thousand},
 pdfsubject={A reproducible arithmetic proof for the negative shift on a ternary sequence}}
\title{Transcendence of shifted Mills minima\\for odd multipliers up to one thousand}
\author{Research manuscript prepared for Keio}
\date{3 November 2026}
\begin{document}
\maketitle
\begin{abstract}
For a positive integer $r$, let $\xi_r$ be the least real $A>1$ for which every
$\lfloor A^{r3^k-1}\rfloor$, $k\ge1$, is prime. We prove that $\xi_r$ is
transcendental for every odd $r\le1000$. The case $r=273$ was expressly
left undecided by the proof in a previous corpus manuscript. Starting from
Saito's Pisot reduction, a local trace identity confines an algebraic minimum
at the exceptional multiples of thirteen to cubics of a specific form.
Primality at an early, prescribed exponent bounds their coefficients.
An exhaustive covering by residue classes then gives a recurring prime
divisor for every surviving cubic. A standalone integer-arithmetic
verification is included.
\end{abstract}

\section{The question and the result}
The starting index matters for a least prime-representing constant. Throughout,
we keep every exponent in the definition
\begin{equation}\label{eq:minimum}
 C_k=r3^k-1\quad(k\ge1),\qquad
 \xi_r=\min\{A>1:\lfloor A^{C_k}\rfloor\text{ is prime for every }k\ge1\}.
\end{equation}
The minimum exists by Saito's Theorem~1.3, since $C_1\ge2$ and the
successive ratios exceed three \cite{saito}. Saito proved transcendence
unconditionally for $r\ge4.003\cdot10^{14}$, and for the whole family
under the Riemann hypothesis. His result identifies a specific number
at each $r$, rather than merely producing some admissible transcendental number.

\begin{theorem}\label{thm:main}
For every odd integer $1\le r\le1000$, the number $\xi_r$ in
\eqref{eq:minimum} is transcendental. In particular, the least $A>1$ for
which $\lfloor A^{273\cdot 3^k-1}\rfloor$ is prime at every $k\ge1$ is
transcendental.
\end{theorem}

Keio's corpus manuscript~\cite{keio64} derives transcendence for odd
$13\nmid r$ and also treats the two fixed shifts $3^k\pm1$.
Manuscript~\cite{keio75} handles the additional multipliers
$13,39,65,91,195$. Its discussion after the $195$ argument singles out
$273$. Although the $273$ sequence is a tail of the $91$ sequence,
removing the first prime constraint can change the minimum.
In particular $9\cdot91-1=818$, the second exponent for $91$, is the first
for $273$, while $3\cdot91-1=272$ has disappeared. The proof below
uses primality at an index actually present for $273$.
The theorem supplies a common verification for all 38 odd multiples of
thirteen in the stated range, 33 beyond the five cases above.

The local return
at three, the cubic descent, the parity restriction, and the height estimate
come from \cite{keio64,keio75}. We prove them here from Saito's results and
supply a uniform exhaustive final step. No implication from transcendence
of either fixed minimum to transcendence of a different minimum is assumed.
All references to Saito's numbering use arXiv version~3.

\section{A cubic from the prime floors}
Suppose $r$ is odd and $A=\xi_r$ is algebraic. Put $B=A^2$ and
$m_k=C_k/2$. Squaring gives an increasing bijection of the respective
admissible sets, so $B$ is the minimum for $(m_k)$. We have
$m_{k+1}=3m_k+1$. For every $i$ some $h>0$ satisfies
$3^h\equiv1\pmod {2m_i}$, since $2m_i=r3^i-1$ is prime to three.
Thus $m_i$ properly divides $m_{i+h}$. An integral $B^{m_i}$ would
make the prime floor at $i+h$ a proper power.

Apply Saito's Proposition~3.1 with his unconditional Theorem~1.6 and
$\theta=21/40$. The exponent ratios exceed three and tend to three,
so all sufficiently large indices belong to his index set. The power
$B^g$ in that proposition has $g\mid m_k$ eventually, hence $g=1$.
Its degree is at most $1+40/17<4$ and is at least two. Degree two would
require all sufficiently late $m_k$ to be odd primes, whereas $m_k$
is even infinitely often. Consequently $B$ is a cubic Pisot number.
Writing its conjugates as $B_1=B,B_2,B_3$, we have
\[
 P_k=\lfloor B^{m_k}\rfloor=\sum_{i=1}^3 B_i^{m_k}
 \qquad(k\text{ sufficiently large}).
\]
Here a Pisot number is a real algebraic integer greater than one whose
other conjugates have modulus less than one. If $B_2,B_3$ were both
real, their contributions at every even exponent would have positive
sum, contrary to the floor identity. Thus $\{B_2,B_3\}=\{Re^{i\theta_0},Re^{-i\theta_0}\}$
with $0<\theta_0<\pi$ and $0<R<1$.

We will need the quantitative bound $R\le B^{-17/40}$. For each
$0<\kappa<17/40$ the fractional-part estimate in Saito's Proposition~3.1(i)
gives
\[
 |2 R^{m_k}\cos(m_k\theta_0)|=O_\kappa(B^{-\kappa m_k}),
 \qquad (1-21/40)\lim m_{k+1}/m_k-1=17/40 .
\]
If $R>B^{-\kappa}$ then $\cos(m_k\theta_0)\to0$. It follows from the
triple-angle formula that $\cos(3m_k\theta_0)\to0$ and
$|\sin(3m_k\theta_0)|\to1$. But
$|\cos(m_{k+1}\theta_0)|=|\cos(3m_k\theta_0+\theta_0)|\to\sin\theta_0>0$.
This proves the bound by letting $\kappa$ increase to $17/40$.

\section{The local identity at three}
Choose square roots $\gamma_i^2=B_i$ in a fixed normal number
field $K$. The eventual, increasing primes are
$P_k=\sum_i\gamma_i^{r3^k-1}$. Discard finitely many $P_k$ that
divide $3\prod_i N_{K/\mathbb Q}(\gamma_i)$. At a prime of
$K$ above $P_k$, let $d\le[K:\mathbb Q]$ be the residue degree and
$E=P_k^d-1$. If $v_3(E)\le k+v_3(r)$, there exists $h>0$ with
$r3^{k+jh}\equiv r3^k\pmod E$ for all $j\ge0$, by taking $3^h$
congruent to one modulo the prime-to-three part of $E$. All
$\gamma_i$ are residue-field units, so $P_k\mid P_{k+jh}$, a
contradiction for large $j$. Thus $v_3(P_k^{2d}-1)>k+v_3(r)$. Since
$v_3(P_k^{2d}-1)=v_3(P_k^2-1)+v_3(d)$ and $d$ is bounded, we get
\[
 \max\{v_3(P_k-1),v_3(P_k+1)\}\ge k-O_K(1).
\]

Fix a place of $K$ above three with residue degree $f$. For its set
$U$ of unit roots $\gamma_i$, let $\omega_i$ be the Teichm\"uller
representative of the residue and put $v_i=\omega_i^r$. Along
$k\equiv j\pmod f$ as $k\to\infty$, the powers of
$\gamma_i/\omega_i$ with exponent $r3^k$ tend to one. Nonunit
terms tend to zero. Since $P_k$ approaches the set $\{-1,1\}$, we obtain
for every $j\ge0$ an equality in algebraic numbers
\begin{equation}\label{eq:d13two}
  E_j=\sum_{i\in U} v_i^{3^j}/\gamma_i\ \in\{-1,1\}.
\end{equation}
The representatives are roots of unity and hence algebraic. Thus the
equalities in the completion are transported through any complex
embedding of the resulting enlarged number field.


The set $U$ cannot be empty. If it is a proper subset, take an
automorphism of a common normal number field sending one of the
$B_i$ indexed by $U$ to a nonunit root. Such an automorphism exists
because the cubic is irreducible. When $|U|=1$ there is just one
transformed term. When $|U|=2$ there is just one nonunit among the
three roots, so exactly one transformed term has negative valuation.
All numerators are roots of unity. In either case the sum has negative
valuation, contrary to \eqref{eq:d13two}. Consequently all roots are
units at three.
\section{The cubic modulo three}
Let $F=X^3-uX^2+vX-w$ be the polynomial of $B$. If $F$ is
reducible modulo three, all its residual roots belong to
$\mathbb F_9^\times$. When they all belong to $\mathbb F_3^\times$,
the numbers $v_i^2$ are signs. If these signs agree, the eventual
traces $\sum B_i^{m_k}$ vanish modulo three, since their three
residues coincide ($r$ is odd). Otherwise one sign occurs once.
One of $E_0+E_1,E_0-E_1$ isolates twice that term, since
$v_i^3=v_i v_i^2$. The isolated term is nonzero, so its doubled value is $\pm2$. This makes one $\gamma_i$ a root of unity.

Otherwise there is one linear root and a conjugate quadratic pair of
order either eight or four. If the quadratic roots have order eight,
their $v_i$ have order sixteen, while the linear $v_i$ has order at
most four. Hence $E_0+E_2$ isolates the linear term, with the same
contradiction since $v_i^9=-v_i$ on the quadratic terms and $v_i^9=v_i$ on the linear term. If the quadratic roots have order four then the two
values $v_i^2$ there are negatives of each other. The determinant of
their columns in the system with rows $E_0,E_1$ has modulus two
in every complex embedding, since the determinant is $v_2v_3(v_3^2-v_2^2)$ after indexing the pair as $2,3$. Choose a complex embedding carrying the $B_i$ above the linear residual root to $B$. Subtracting its term from each row leaves right sides of modulus at most $1+1/A$. Cramer's rule
bounds each of the two other $|1/\gamma_i|$ by $1+1/A$.
Thus $A<3$, since each small $|1/\gamma_i|$ is at least $A^{17/40}$
and $3^{17/40}>4/3$. Consequently $w=N(B)=1$ since
$0<w\le B^{3/20}<9^{3/20}<2$. The small conjugates
of $B$ now have modulus $B^{-1/2}$, so $\sqrt A\le1+1/A$, which implies $A<\sqrt5$ (at $A=\sqrt5$, $5^{1/4}>1+1/\sqrt5$).
Consequently $B<5$. As $w=1$, the reduction of $F$ modulo three equals
$(X-1)(X^2+1)$ because the quadratic roots multiply to one and the
product of all roots is one. Hence $u\equiv v\equiv1\pmod3$.
The elementary symmetric sums give $-1<u<7$ and
$|v|<2\sqrt5+1<6$. Among $u\in\{1,4\}$ and
$v\in\{-5,-2,1,4\}$, the tests $F(1)<0<F(5)$ and
$\operatorname{disc}(F)<0$ leave precisely
$(u,v)=(1,-2),(4,-2),(4,1)$.

For these three cubics let $T_n=\sum_i B_i^n$. We eliminate each at modulus $27$. In $(\mathbb Z/27)[X]/F$, the
class $x$ is a unit. Modulo three $F$ is squarefree
and its roots belong to $\mathbb F_9^\times$, so $x^8=1$ modulo
three, and two cubings suffice to lift this to $x^{72}=1$ modulo
27. In particular the claimed exponent $216$ is valid. Thus the
traces $T_{m_k}$ modulo 27 for all sufficiently large $k$ depend
only on $r\bmod16$ and $k\bmod4$, since $216=8\cdot27$ and one may
compute $m_k$ using $r3^k-1\bmod432$. The following standard-library
calculation checks that some recurrent phase is neither sign in each case.
\begin{lstlisting}[language=Python]
for u,v in ((1,-2),(4,-2),(4,1)):
    T=[3,u,(u*u-2*v)%27]
    while len(T)<216:
        T.append((u*T[-1]-v*T[-2]+T[-3])%27)
    for r0 in range(1,16,2):
        residues=[T[((r0*3**k-1)//2)%216] for k in range(6,10)]
        assert any(t not in (1,26) for t in residues)
\end{lstlisting}
This contradicts the congruence for $P_k$, proving that $F$ is
irreducible modulo three.

\section{The two cases at thirteen}\paragraph{Case $13\nmid r$.}
Write $\mu_n$ for the group of $n$th roots of unity. The three residual $B_i$ form one Frobenius orbit in $\mathbb F_{27}^{\times}$. If $13\nmid r$,
this leads to a contradiction. Write the Teichm\"uller lifts for
$\gamma_i$ as $\omega_i=s_i t_i$, with $s_i\in\mu_4$ and
$t_i\in\mu_{13}$. The squares $s_i^2$ are equal, since they are
the sign components of the Frobenius-conjugate lifts of $B_i$.
Changing the signs of $\gamma_i$ when necessary makes all $s_i$ equal
to a common $s$. This is legitimate because every $r3^k-1$ is even.
The $13$-components $t_i^2$ form a Frobenius orbit, and squaring is
a bijection on $\mu_{13}$. Thus, after indexing, the $t_i$ are
$t,t^3,t^9$, where $t$ is primitive because the residual $B_i$
have degree three.

Divide equation~\eqref{eq:d13two} at $j=0,1,2$ by $s^{r3^j}$.
The coefficient rows for the three unknowns $1/\gamma_i$ are the
three cyclic shifts of $(\eta,\eta^3,\eta^9)$, with $\eta=t^r$
primitive. To check that all its Fourier factors are nonzero, take
$w\in\mu_3$. The expression $\eta+w\eta^3+w^2\eta^9$ is the
evaluation at $\eta$ of a nonzero polynomial of degree nine over
$\mathbb Q(\mu_3)$, whereas the relative degree of $\eta$ is
$\varphi(39)/\varphi(3)=12$. The right sides and all coefficients
belong to $\mathbb Q(\mu_{52})$, so the solution does too. In
particular the real cubic field $\mathbb Q(B)$ lies in an
abelian extension of $\mathbb Q$ and is Galois. As a real Galois
field it is totally real. This contradicts the nonreal conjugate
pair established above. Consequently the minimum is transcendental
for every odd $r$ with $13\nmid r$.
\paragraph{Case $13\mid r$.}
The residual roots of $B$ are now a Frobenius orbit in
$\mathbb F_{27}^\times$, of order dividing $26$. Because $13\mid r$,
the three $v_i^2$ are equal to a common sign $s$ (all are signs,
and Frobenius permutes them). Put $y_i=v_i/\gamma_i$. By
\eqref{eq:d13two}, $\sum y_i=\pm1$ and $y_i^2=s/B_i$.
Any Galois automorphism of a normal field containing these numbers
sends each $y_i$ to some $\pm y_{\pi(i)}$. A sign change would give
a nonempty vanishing subsum of $\sum y_i$, by subtracting the two
conjugate equalities. Neither one nor three terms can vanish, and
two cannot vanish since the $B_i$ are distinct. Thus there are
no sign changes. In particular $y_1\in\mathbb Q(B)$, which is
real at the embedding with $B>1$, giving $s=1$.
Consequently $A=\pm y_1^{-1}$ is cubic Pisot and
$\operatorname{Tr}(A^{-1})=\pm1$. Write its minimal polynomial as
$X^3-aX^2+bX-c$. The trace of the inverse is $b/c$, so $b=\pm c$.
The conjugate bound gives $|N(A)|\le A^{3/20}$ and its nonreal
conjugate pair gives positive norm. Finally the residues of the
three $B_i$ are distinct of degree three, so the residue of
$A$ also has degree three (as $A^2=B$). Its own minimal
polynomial is therefore irreducible over $\mathbb F_3$. Traces of
the powers along the unchanged clock agree eventually with the prime floors because $A^2=B$ and both generate the same cubic field.
\section{A coefficient bound from an early prime floor}
Suppose now that $13\mid r$. Write $F(X)=X^3-aX^2+bX-c$
for the polynomial of the minimum $A$. Write $T_m$ for its trace at exponent $m$.
Every trace $T_{r3^k-1}$, including the early ones, is odd. Indeed,
eventually they are odd primes, and $b=\pm c$. Modulo two, if $c$ is even
the trace at $m\ge1$ is $a^m$. If $a,c$ are both odd, $F=(X+1)^3$
in characteristic two and all these traces equal one. If $a$ is even
and $c$ odd, $F=X^3+X+1$ modulo two and the roots have order seven.
Then $T_{r3^k-1}$ modulo two is periodic from the first index because
$3$ is a unit modulo seven. The same observation shows that $a$ is
odd when $7\nmid r$. Indeed, if $a$ is even we must have $c$ odd.
Since $3$ generates $(\mathbb Z/7)^\times$, there is a phase with
$r3^k-1\equiv1\pmod7$, giving the even parity of $a$ and a contradiction.

Set $N=3r-1$ if $r\equiv3\pmod4$, and $N=9r-1$ otherwise, so $4\mid N$.
Write the small conjugates of $A$ as $\sqrt q e^{\pm i\theta}$, with
$q=c/A\le A^{-17/20}$ and $0<\theta<\pi$. The trace of $1/A$ is
$\pm1$, whence $\cos\theta=\sqrt q(\pm1-1/A)/2$.
If $A\ge(N/2)^{40/17}$, then $A>4$ and
$t=|\cos\theta|<5/(4N)<1/2$. Thus
$N|\theta-\pi/2|=N\arcsin(t)<6Nt/5<3/2$.
As $N$ is divisible by four, $\cos(N\theta)>0$, and
\[
 0< T_N-A^N=2q^{N/2}\cos(N\theta)\le2(2/N)^N<1.
\]
Therefore $\lfloor A^N\rfloor=T_N-1$ is even and larger than two,
contradicting its required primality at the original index. We conclude
\[
 A<(N/2)^{40/17},\qquad c<(N/2)^{6/17}.
\]
This argument uses no floor-to-trace identification at early indices.

The sum of the three conjugates gives $A-2<a<A+2$. Since $A>1$,
the possible integer values lie between zero and
$\lfloor(N/2)^{40/17}\rfloor+2$, inclusive. We may also include
$c=\lfloor(N/2)^{6/17}\rfloor$ when forming an outer bound. The
conditions $3\nmid c$, $b=\pm c$, irreducibility modulo three, and the
parity restrictions above define the finite grid used below.
\section{A finite covering of cubic coefficients}
Let $f=X^3-aX^2+bX-c$ be a cubic in the grid. For a prime $p\ge5$ with $p\nmid c$, the class $x$
of $X$ is invertible in $R=\mathbb F_p[X]/(f)$ and satisfies
$x^{E_p}=1$, where
\[
 E_p=p(p^2-1)(p^2+p+1).
\]
In the reduced quotient each residue-field multiplicative group has
order $p^j-1$ for $1\le j\le3$, all dividing $E_p/p$.
The nilradical $J$ satisfies $J^3=0$, so $u^p=1$ for every $u\in1+J$.
Let $e=v_3(E_p)$. If $k\ge e$ and
$\operatorname{Tr}_R(x^{r3^k-1})=0$, take any $h>0$ with
$3^h\equiv1\pmod{E_p/3^e}$. The difference of exponents at
$k+jh$ and $k$ is $r3^k(3^{jh}-1)$, divisible by $E_p$. Therefore
$p\mid T_{r3^{k+jh}-1}$ for all $j\ge0$. Since $A$ is Pisot, the
traces grow to infinity. A single such certificate excludes the entire
coefficient residue class $a\pmod p$ for these $b,c$.

The checker in Appendix~\ref{app:code} searches $p\le1000$ and
$\max(1,e)\le k\le100$. It performs polynomial multiplication in the
basis $1,x,x^2$, computes the algebra trace from
$\operatorname{Tr}_R(1)=3$, $\operatorname{Tr}_R(x)=a$, and
$\operatorname{Tr}_R(x^2)=a^2-2b$, and uses
$x^{r3^{k+1}-1}=x^2(x^{r3^k-1})^3$.
Primes dividing $c$ are always skipped. For each fixed pair
$(b,c)$ the residue classes obtained in this way are removed from the
grid. The surviving count is an exact integer calculation. An independent check
used powers of the $3\times3$ companion matrix modulo each witness prime
and separately covered the coefficient intervals. Both calculations
gave the same counts for all 38 multipliers.

\section{What one residue class certifies}
For $r=273$ the first exponent divisible by four is $N=9r-1=2456$.
The grid has $a\le18\,565\,304$, $c\le12$, and $3\nmid c$.
As $7\mid273$, even $a$ must be retained when $c$ is odd.
Consider, for example, $f=X^3-X^2-X-1$. At $p=5$ and $k=6$ its exponent is
$273\cdot3^6-1=199\,016$. Reduction by $x^3=x^2+x+1$ gives
\[
 x^{199\,016}=2+3x+2x^2\quad\text{in }\mathbb F_5[X]/(f).
\]
Its trace is $2\cdot3+3\cdot1+2\cdot3=15\equiv0\pmod5$.
Here $E_5=3720$, $v_3(E_5)=1$, and $3^{60}\equiv1\pmod{1240}$.
The same prime divides every trace at $k=6+60j$, $j\ge0$. This rules
out the whole class $a\equiv1\pmod5$ for $b=-1$, $c=1$.
The computation in the appendix performs the same operation with the
witness prime and index chosen separately for each class.

\begin{center}
\begin{tabular}{lrrr}\toprule
Multiplier(s) & Triples in the grid & Rules used & Largest prime\\\midrule
$13$ & $21\,153$ & $388$ & $53$\\
$117$ & $7\,576\,020$ & $2\,562$ & $103$\\
$273$ & $111\,391\,828$ & $4\,188$ & $179$\\
$949$ & $2\,265\,148\,054$ & $11\,577$ & $173$\\\midrule
All 38 & $13\,197\,086\,314$ & $181\,239$ & $241$\\\bottomrule
\end{tabular}
\end{center}
Counts are per multiplier and then summed. In every row the number of
uncovered triples is zero. Only eventual equality of floor and trace
enters the exclusion, even when the first residue witness occurs early.
Together with the case $13\nmid r$ this proves Theorem~\ref{thm:main}.

The size limit $1000$ specifies the range of the exhaustive verification.
For a further odd multiple of thirteen the reduction and the height bound
still apply, but the resulting finite grid needs its own covering.
Even multipliers, including $r=2$, require a different reduction at
the outset because $C_k$ is odd. The classical Mills minimum has the
different exponent sequence $(3^k)_{k\ge1}$.

\begin{thebibliography}{9}
\bibitem{saito} K. Saito, \emph{Transcendency of variants of Mills' constant},
Ramanujan J. \textbf{70} (2026), article 69.
\href{https://doi.org/10.1007/s11139-026-01443-0}{doi:10.1007/s11139-026-01443-0}.
Theorem numbering in this paper refers to
\href{https://arxiv.org/html/2508.16068v3}{arXiv:2508.16068v3}.
\bibitem{keio64} Keio, \emph{Ternary trace dynamics and the transcendence of shifted
Mills constants}, Verifiable Mathematical Reasoning Corpus, manuscript 064,
25 October 2026, Theorem B and Section 7.2.
\bibitem{keio75} Keio, \emph{Shifted Mills Constants and Cubic Pisot Trace Dynamics},
Verifiable Mathematical Reasoning Corpus, manuscript 075, 27 October 2026,
in particular the discussion of $r=273$ following the proof for $r=195$.
Both corpus references use the
\href{https://github.com/QEEDDD/Verifiable-Mathematical-Reasoning-Corpus/tree/ecd87de3348dcb706a76bb12fde722d7a348a012}{2 November 2026 repository revision}.
\end{thebibliography}

\appendix
\section{Reproducing the finite covering}\label{app:code}
The following C++17 program uses only the standard library. The five
output columns are $r$, grid size, number of covering rules, largest
witness prime, and number left. Pass
\texttt{273} to the executable for the single case or pass no argument
for all 38. For instance, the following commands extract the listing
from this TeX source and execute it.
\begin{verbatim}
awk '/^%% BEGIN CHECKER$/ {f=1;next} /^%% END CHECKER$/ {f=0}
     f && $0 !~ /^\\(begin|end)\{lstlisting\}/ {print}' \
     Shifted_Mills_Minima.tex > check.cpp
c++ -O2 -std=c++17 check.cpp -o check
./check 273
./check
\end{verbatim}
The covering loop keeps a list of progressions modulo a common width.
At a new prime $p$, each progression splits into at most $p$ classes.
It retains only those whose trace has no witness at this prime. Once
the width exceeds the coefficient bound, each class is a singleton.
This explains why the memory cost is smaller than the initial grid size.
The two integer root bounds are computed by comparing exact powers in
radix $10^8$. The anticipated final line is
\begin{verbatim}
total 13197086314 181239 241 0
\end{verbatim}

%% BEGIN CHECKER
\begin{lstlisting}[language=C++]
// C++17. With no argument, checks all 38 multipliers. With e.g. 273, checks one.
#include <algorithm>
#include <array>
#include <cstdint>
#include <cstdlib>
#include <iostream>
#include <vector>

using Int = std::int64_t;
using Poly = std::array<int,3>;

std::vector<int> exact_power(int base, int exponent) {
    // The digit count followed by the most significant digits compares lexically.
    constexpr int radix = 100000000;
    std::vector<int> digits{1};
    for (int i=0; i<exponent; ++i) {
        Int carry=0;
        for (int& d: digits) {
            carry += Int(d)*base;
            d = int(carry%radix);
            carry /= radix;
        }
        while (carry) {
            digits.push_back(int(carry%radix));
            carry /= radix;
        }
    }
    digits.push_back(int(digits.size()));
    std::reverse(digits.begin(),digits.end());
    return digits;
}

int bound(int base, int power) {
    auto target=exact_power(base,power);
    int lo=0, hi=1;
    while (exact_power(hi,17)<=target) hi*=2;
    while (hi-lo>1) {
        int mid=lo+(hi-lo)/2;
        if (exact_power(mid,17)<=target) lo=mid;
        else hi=mid;
    }
    return lo;
}

bool prime(int n) {
    for (int d=2; d*d<=n; ++d) if (n%d==0) return false;
    return n>=2;
}

struct Cubic {
    int a,b,c,p;
    int mod(Int n) const { return int((n%p+p)%p); }
    Poly multiply(Poly u, Poly v) const {
        // Reduce the degree-four product using X^3=aX^2-bX+c.
        Int h4=Int(u[2])*v[2];
        Int h3=mod(Int(u[1])*v[2]+Int(u[2])*v[1]+a*h4);
        return {mod(Int(u[0])*v[0]+c*h3),
                mod(Int(u[0])*v[1]+Int(u[1])*v[0]-b*h3+c*h4),
                mod(Int(u[0])*v[2]+Int(u[1])*v[1]+Int(u[2])*v[0]+a*h3-b*h4)};
    }
    int witness(int r, int depth) const {
        const Poly x2{0,0,1};
        Poly x{0,1,0}, u{1,0,0};
        for (int n=3*r-1; n; n/=2) {
            if (n%2) u=multiply(u,x);
            x=multiply(x,x);
        }
        for (int k=1; k<=100; ++k) {
            Int trace=3*u[0]+Int(a)*u[1]+Int(a*a-2*b)*u[2];
            if (k>=depth && mod(trace)==0) return k;
            u=multiply(x2,multiply(u,multiply(u,u)));
        }
        return 0;
    }
};

struct Counts { Int grid=0, rules=0, left=0; int largest=0; };

void check_row(int r,int b,int c,int A,Counts& counts) {
    // A representative j denotes all j+width*t in [0,A], t>=0.
    std::vector<int> live;
    Int width=6;
    for (int d=0; d<6; ++d) {
        if (d%2==0 && (r%7!=0 || c%2==0)) continue;
        bool root=false;
        for (int x=0; x<3; ++x)
            if ((x*x*x-d*x*x+b*x-c)%3==0) root=true;
        if (!root) { live.push_back(d); counts.grid+=(A-d)/6+1; }
    }
    for (int p=5; p<=1000 && !live.empty(); p+=2) {
        if (!prime(p) || c%p==0) continue;
        int depth=(p%3==1);
        for (int t=p*p-1; t%3==0; t/=3) ++depth;
        std::vector<int> cache(p,-1), next;
        for (int j:live) {
            Int last=std::min<Int>(A,j+(p-1)*width);
            for (Int a=j; a<=last; a+=width) {
                int v=int(a%p);
                if (cache[v]<0) {
                    cache[v]=Cubic{v,b,c,p}.witness(r,depth);
                    if (cache[v]) {
                        ++counts.rules;
                        counts.largest=std::max(counts.largest,p);
                    }
                }
                if (cache[v]==0) next.push_back(int(a));
            }
        }
        width=std::min<Int>(Int(A)+1,width*p);
        live.swap(next);
    }
    for (int j:live) counts.left+=(A-j)/width+1;
}

int main(int argc,char** argv) {
    int only=argc==2?std::atoi(argv[1]):0;
    if (argc>2 || (argc==2 && (only<13 || only>1000 || only%26!=13))) return 2;
    Counts all;
    for (int r=13; r<=1000; r+=26) {
        if (only && only!=r) continue;
        int N=(r%4==3?3:9)*r-1;
        int A=bound(N/2,40)+2, C=bound(N/2,6);
        Counts row;
        for (int c=1; c<=C; ++c) if (c%3)
            for (int b:{-c,c}) check_row(r,b,c,A,row);
        std::cout<<r<<" "<<row.grid<<" "<<row.rules<<" "
                 <<row.largest<<" "<<row.left<<std::endl;
        all.grid+=row.grid; all.rules+=row.rules; all.left+=row.left;
        all.largest=std::max(all.largest,row.largest);
    }
    std::cout<<"total "<<all.grid<<" "<<all.rules<<" "
             <<all.largest<<" "<<all.left<<std::endl;
    return all.left!=0;
}
\end{lstlisting}
%% END CHECKER
\end{document}
